% =====================================================================
% Response --- Reviewer 1
% Source: docs/reviews/metro-rev1.docx
% Point map: docs/DOCS_ASSESSMENT.md, Stage 3 (R1.1-R1.9).
% Status: STUB --- points restated, responses to be written.
% =====================================================================
\section*{Reviewer 1}

\noindent General assessment: the paper is interesting and of interest to
Metroeconomica readers, but its structure does not let the reader focus on
the primary message; the central argument is not yet discernible.

\subsection*{R1.1 --- Structure and central argument}
\refcom{I find the paper challenging to follow, and its current structure
does not facilitate the reader's focus on the primary message.}
\response TODO: point to the new spine (Sec. 1 intro + Sec. 5 matrix) and
the reading-guide box.
\loc Sec. 1; Sec. 5.

\subsection*{R1.2 --- Verbal description of the model; static vs.\ dynamic}
\refcom{It would be nice to say verbally what are the main features of the
model \dots\ The default model is a closed economy with full employment
\dots\ is the model static or dynamic?}
\response TODO: verbal model overview before the equations; state the static
short-run horizon explicitly.
\loc Sec. 4.1.

\subsection*{R1.3 --- Calibration process and summary tables}
\refcom{How the model has been calibrated? Is the model exactly replicating
the IO table? \dots\ present key economic variables from Germany's IO table
in a summary table \dots\ Including elasticity values in a table.}
\response TODO: the calibration pipeline (Sec. 4.1); the summary table; the
elasticity table.
\loc Sec. 3 / App. B; elasticity table.

\subsection*{R1.4 --- The Keynesian approach (Type I vs Type II multipliers)}
\refcom{The purpose behind the Keynesian approach in section 2.1 remains
unclear \dots\ distinguishing between Type I and Type II multipliers
(Miller and Blair, 2019).}
\response TODO: clarify or drop the Type-I/II distinction.
\loc Sec. 2 / Sec. 3.

\subsection*{R1.5 --- Labour-market slack and the wage-setting rule}
\refcom{This model evidently allows for involuntary unemployment \dots\ a
wage-setting rule must be specified \dots\ some wage-setting mechanism such
as fixed real or nominal wages \dots\ Also a simple wage curve.}
\response TODO: GAMMA/BETA with explicit wage rules; the no-go corollary on
nominal wage rules.
\loc Sec. 4 (closures); Sec. 6.

\subsection*{R1.6 --- Labour mobility and the unified wage rule}
\refcom{Under imperfect labour market conditions, does the model assume a
unified wage rule to enable labour mobility across sectors?}
\response TODO: BF/ALPHA; the sectoral family (ADR-0022).
\loc Sec. 4; Sec. 6.

\subsection*{R1.7 --- Aggregate vs.\ sectoral results}
\refcom{In aggregate the economy is not sensitive to change in elasticity of
substitution while these changes in price sensitiveness matter a lot for
sectoral outcome.}
\response TODO: sections 5-6 with the sectoral Sobol.
\loc Sec. 5.1-5.2; Sec. 6.

\subsection*{R1.8 --- McGregor, Swales \& Yin (1996)}
\refcom{[the long-run interpretation of regional IO analysis]}
\response TODO: engage; the DELTA corner and the long-run bridge.
\loc Sec. 2 (literature); Sec. 6.

\subsection*{R1.9 --- Minor comments}
\refcom{Page 5 sentence is convoluted; Page 16 second paragraph incomplete;
Figure 3 --- is the horizontal axis representing sectors?}
\response TODO: copy-edits.
\loc Sec. 2; Sec. 4; Figure (sectors on the axis).